\chapter{Conclusion}
\label{chap:conclusion}



\section{Conclusion}

This study project investigated whether anomaly detection can operate
without explicit protocol knowledge while maintaining strong
performance on reverse-engineered traffic representations.

The results show that protocol-agnostic detection is feasible,
particularly when operating on raw packet data.
Autoencoder-based detection, CNN models, and N-gram approaches
achieve consistently strong performance, especially in
Scenario~2 where sufficient attack diversity is available.

Reverse engineering reduces computational cost but introduces
representation-dependent performance trade-offs.
While deep learning models remain comparatively robust under
moderate compression, traditional machine learning classifiers
experience significant degradation.
The most aggressively compressed representation (RE15) yields
the weakest overall results, indicating that critical structural
information is lost.

Scenario characteristics further influence performance:
limited attack diversity (Scenario~3) in the training data consistently reduces
generalization capability across all model families.

In summary, raw packet representations provide the most reliable
detection performance, whereas reverse-engineered representations
offer efficiency gains at the cost of reduced robustness.
The optimal choice therefore depends on the trade-off between
computational constraints and detection accuracy.

\section{Limitations}

Several limitations should be considered when interpreting the
results of this study project.

First, a post-hoc analysis revealed structural characteristics of
the dataset that may have influenced detection performance.
In particular, the attacker IP address appears to remain constant
across attack samples. Since the RAW representation includes full
packet headers, including IP-level information, models trained on
RAW data may have partially exploited this invariant identifier.
In contrast, the reverse-engineered representations exclude such
header information, which may contribute to the observed performance
gap between RAW and RE-based datasets. Consequently, the absolute
performance of RAW-based models should be interpreted with caution.

Second, attack labels were derived from the dataset folder structure
rather than from packet-level ground-truth annotations. As a result,
benign server responses that occurred within attack sessions were
also assigned the attack label. This introduces label noise and may
have affected both training and evaluation, particularly for models
sensitive to fine-grained traffic structure.

{\color{gray} Third, the reverse-engineered representation was designed to extract
a dynamic region within S7Comm payloads that was assumed to contain
process-related information. However, since the dataset does not
provide explicit ground-truth annotations for physical sensor values,
it cannot be guaranteed that the extracted region consistently
corresponds to actual physical readings. The extracted sequences
should therefore be interpreted as reverse-engineered payload
segments rather than verified process measurements.}
{\color{red} todo: describe physical readings problem correctly }


Finally, the evaluation was conducted on a specific dataset and
predefined scenarios. The observed trends may vary under different
network environments, attacker diversity, or alternative
reverse-engineering strategies. Future work should incorporate
datasets with greater attacker variability, stricter packet-level
labeling, and validated process-value annotations to further
assess robustness.

\section{Future Work}

The results of this study project open several avenues for further investigation.

The reverse-engineered representations rely on byte-level concatenation
of extracted payload segments. Although aggregation controlled by the
parameter $p$ increases contextual information, it may obscure
fine-grained temporal and structural dependencies within individual
packets. Future work could explore sequence-aware modeling approaches
that preserve packet boundaries or explicitly encode temporal order
instead of relying solely on concatenation.

The traditional ML classifiers operate on latent features extracted
by the stacked denoising autoencoder. While this compact representation
proved effective in certain settings, alternative feature extraction
strategies may further improve performance. For example, statistical
traffic descriptors, protocol-agnostic sequence embeddings, or
transformer-based encoders could provide richer and more expressive
representations for classical machine learning models.

Moreover, the evaluation was conducted on a dataset with specific
structural characteristics, including limited attacker diversity and
coarse-grained labeling. Future studies should therefore consider
more heterogeneous datasets with diversified attacker behavior,
strict packet-level annotations, and validated process-value ground
truth. Such validation would allow for a more robust assessment of
generalization performance under realistic industrial conditions.